% ============================================================
% Close-up: OmicMiningAgent — Workflow Interface View
% Fragment — % ============================================================
% Close-up: OmicMiningAgent — Workflow Interface View
% Fragment — % ============================================================
% Close-up: OmicMiningAgent — Workflow Interface View
% Fragment — % ============================================================
% Close-up: OmicMiningAgent — Workflow Interface View
% Fragment — \input{figs/closeup_omic}
% Complements internal tool illustrations with orchestrator-
% level data flow, prompt highlights, and state I/O.
% ============================================================
\begin{figure}[htbp]
\centering
\resizebox{\textwidth}{!}{%
\begin{tikzpicture}[node distance=0.35cm, font=\sffamily,
    stepbar/.style={rectangle, rounded corners=6pt, thick,
        minimum width=19cm, align=center, font=\sffamily,
        drop shadow={opacity=0.10}, inner sep=6pt, text width=18cm},
    rwpanel/.style={rectangle, rounded corners=5pt, thick,
        minimum width=19cm, align=left, inner sep=6pt,
        font=\sffamily\small, text width=17.5cm},
    infopanel/.style={rectangle, rounded corners=5pt, draw=gray!50, thick,
        minimum width=19cm, align=left, fill=white,
        inner sep=6pt, font=\sffamily\small, text width=17.5cm},
    toolpanel/.style={rectangle, rounded corners=5pt, draw=gray!60, thick,
        minimum width=19cm, align=left, fill=gray!4,
        inner sep=6pt, font=\sffamily\small, text width=17.5cm},
    flowline/.style={->, >=Stealth, very thick, draw=gray!40},
    consumerlabel/.style={font=\sffamily\small\itshape, text=gray!60!black,
        align=center, text width=19cm}
]
    % === STEP BAR ===
    \node[stepbar, draw=orange!50!black, fill=agentOrange] (step)
        {\iconDNA{orange!60!black}\;\;
         \textbf{\large OmicMiningAgent --- Data Foundation}
         {\normalsize\itshape\; (typically first)}\\[3pt]
         {\small Planner: \textit{``ALWAYS run FIRST to ground analysis
          in actual expression data --- establishes data foundation.''}}\\[2pt]
         {\scriptsize\sffamily\color{gray!60!black}
          Note: ordering shown is a \emph{typical} realisation;
          the LLM planner may reorder, skip, or loop back to any
          agent based on intermediate results.}};

    % === READS ===
    \node[rwpanel, draw=green!40!black, fill=agentGreen,
          below=0.4cm of step] (reads)
        {\textbf{$\blacktriangledown$ Reads from AgentState}\\[5pt]
         $\bullet$\; \texttt{query} {\scriptsize (str)} ---
          user's natural-language research question\\[2pt]
         $\bullet$\; \texttt{session\_id} {\scriptsize (str)} ---
          unique session identifier for file paths and result caching};

    \draw[flowline] (step.south) -- (reads.north);

    % === PROMPT HIGHLIGHTS ===
    \node[infopanel, below=0.4cm of reads] (prompt)
        {\textbf{System Prompt Highlights}
         {\scriptsize (from \texttt{OmicMiningAgent.system\_message})}\\[5pt]
         \textbf{1.}\; Extract \texttt{disease}, \texttt{organ},
          \texttt{cell\_type} from the user query\\[2pt]
         \textbf{2.}\; Map free-text to canonical forms:\;
          ``Alzheimer's'' $\rightarrow$ \texttt{"Alzheimer disease"},\;
          ``Lung cancer'' $\rightarrow$ \texttt{"lung adenocarcinoma"}\\[2pt]
         \textbf{3.}\; Map organs:\;
          Alzheimer's $\rightarrow$ \texttt{"brain"},\;
          Breast cancer $\rightarrow$ \texttt{"breast"}\\[2pt]
         \textbf{4.}\; After tool calls $\rightarrow$
          ``write a concise summary of the gene list and pathway data''};

    \draw[flowline] (reads.south) -- (prompt.north);

    % === TOOL ===
    \node[toolpanel, below=0.4cm of prompt] (tool)
        {\textbf{Tool:}\;
         \texttt{omic\_analysis\_tool(disease, organ, cell\_type)}\\[5pt]
         {\scriptsize \textbf{Pipeline:}\;
          GLiNER NER $\rightarrow$
          OmniCellTOSG scRNA-seq query $\rightarrow$
          DESeq2 differential expression $\rightarrow$
          clusterProfiler KEGG/GO enrichment}\\[3pt]
         {\scriptsize \textbf{Returns:}\;
          sample counts (disease vs.\ normal), DEG rankings with
          log2FC \& FDR, volcano plots, KEGG/GO enrichment dotplots}\\[3pt]
         {\scriptsize \textbf{Resources:}\;
          OmniCellTOSG\; $\cdot$\;
          GLiNER\; $\cdot$\;
          DESeq2/R\; $\cdot$\;
          clusterProfiler\; $\cdot$\;
          BiomedBERT}};

    \draw[flowline] (prompt.south) -- (tool.north);

    % === WRITES ===
    \node[rwpanel, draw=blue!30!black, fill=stateBlue,
          below=0.4cm of tool] (writes)
        {\textbf{$\blacktriangle$ Writes to AgentState}
         {\scriptsize (via \texttt{\_extract\_structured\_data})}\\[5pt]
         $\bullet$\; \texttt{top\_genes[]} {\scriptsize (list[str])} ---
          ranked DEGs with log2FC \& FDR\quad
          {\scriptsize e.g.\ \texttt{[MAP3K15, CFAP47, IL3RA, \ldots]}}\\[2pt]
         $\bullet$\; \texttt{disease} {\scriptsize (str)} ---
          canonical disease name\quad
          {\scriptsize e.g.\ \texttt{"Alzheimer disease"}}\\[2pt]
         $\bullet$\; \texttt{cell\_type} {\scriptsize (str)} ---
          cell-type filter if specified\quad
          {\scriptsize e.g.\ \texttt{"microglial cell"}}\\[2pt]
         $\bullet$\; \texttt{pathways[]} {\scriptsize (list)} ---
          KEGG and GO enrichment results\\[2pt]
         $\bullet$\; \texttt{deg\_results} {\scriptsize (dict)} ---
          full differential-expression statistics per gene\\[2pt]
         $\bullet$\; \texttt{plot\_paths} {\scriptsize (dict)} ---
          volcano plot, KEGG dotplot, enrichment plot file paths};

    \draw[flowline] (tool.south) -- (writes.north);

    % === DOWNSTREAM CONSUMERS ===
    \node[consumerlabel, below=0.3cm of writes] (consumers)
        {$\downarrow$\; Consumed by:\;
         \textbf{BioMarkerKGAgent} (\texttt{top\_genes})\; $\cdot$\;
         \textbf{PubMedResearcher} (\texttt{top\_genes})\; $\cdot$\;
         \textbf{ScientistsAgent} (\texttt{top\_genes}, \texttt{pathways})\; $\cdot$\;
         \textbf{ReportingNode} (\texttt{plot\_paths})};

\end{tikzpicture}%
}% end resizebox
\caption{\textbf{OmicMiningAgent} --- workflow interface view.
The ordering shown here reflects a \emph{typical} plan generated by
the LLM orchestrator; in practice the planner dynamically decides
which agents to invoke, in what order, and whether to loop back after
failures or surprising results.  That said, the planning prompt
strongly encourages running this agent first to ``ground ALL
hypotheses in actual omics data.''  It reads only the user query and session identifier,
maps them to canonical disease/organ/cell-type forms via its system
prompt, and invokes a single tool (\texttt{omic\_analysis\_tool}) that
chains NER, database retrieval, differential expression, and pathway
enrichment.  The six fields it writes to \texttt{AgentState}---most
critically \texttt{top\_genes[]}---become the shared anchor that
every downstream agent reads.}
\label{fig:closeup-omic}
\end{figure}

% Complements internal tool illustrations with orchestrator-
% level data flow, prompt highlights, and state I/O.
% ============================================================
\begin{figure}[htbp]
\centering
\resizebox{\textwidth}{!}{%
\begin{tikzpicture}[node distance=0.35cm, font=\sffamily,
    stepbar/.style={rectangle, rounded corners=6pt, thick,
        minimum width=19cm, align=center, font=\sffamily,
        drop shadow={opacity=0.10}, inner sep=6pt, text width=18cm},
    rwpanel/.style={rectangle, rounded corners=5pt, thick,
        minimum width=19cm, align=left, inner sep=6pt,
        font=\sffamily\small, text width=17.5cm},
    infopanel/.style={rectangle, rounded corners=5pt, draw=gray!50, thick,
        minimum width=19cm, align=left, fill=white,
        inner sep=6pt, font=\sffamily\small, text width=17.5cm},
    toolpanel/.style={rectangle, rounded corners=5pt, draw=gray!60, thick,
        minimum width=19cm, align=left, fill=gray!4,
        inner sep=6pt, font=\sffamily\small, text width=17.5cm},
    flowline/.style={->, >=Stealth, very thick, draw=gray!40},
    consumerlabel/.style={font=\sffamily\small\itshape, text=gray!60!black,
        align=center, text width=19cm}
]
    % === STEP BAR ===
    \node[stepbar, draw=orange!50!black, fill=agentOrange] (step)
        {\iconDNA{orange!60!black}\;\;
         \textbf{\large OmicMiningAgent --- Data Foundation}
         {\normalsize\itshape\; (typically first)}\\[3pt]
         {\small Planner: \textit{``ALWAYS run FIRST to ground analysis
          in actual expression data --- establishes data foundation.''}}\\[2pt]
         {\scriptsize\sffamily\color{gray!60!black}
          Note: ordering shown is a \emph{typical} realisation;
          the LLM planner may reorder, skip, or loop back to any
          agent based on intermediate results.}};

    % === READS ===
    \node[rwpanel, draw=green!40!black, fill=agentGreen,
          below=0.4cm of step] (reads)
        {\textbf{$\blacktriangledown$ Reads from AgentState}\\[5pt]
         $\bullet$\; \texttt{query} {\scriptsize (str)} ---
          user's natural-language research question\\[2pt]
         $\bullet$\; \texttt{session\_id} {\scriptsize (str)} ---
          unique session identifier for file paths and result caching};

    \draw[flowline] (step.south) -- (reads.north);

    % === PROMPT HIGHLIGHTS ===
    \node[infopanel, below=0.4cm of reads] (prompt)
        {\textbf{System Prompt Highlights}
         {\scriptsize (from \texttt{OmicMiningAgent.system\_message})}\\[5pt]
         \textbf{1.}\; Extract \texttt{disease}, \texttt{organ},
          \texttt{cell\_type} from the user query\\[2pt]
         \textbf{2.}\; Map free-text to canonical forms:\;
          ``Alzheimer's'' $\rightarrow$ \texttt{"Alzheimer disease"},\;
          ``Lung cancer'' $\rightarrow$ \texttt{"lung adenocarcinoma"}\\[2pt]
         \textbf{3.}\; Map organs:\;
          Alzheimer's $\rightarrow$ \texttt{"brain"},\;
          Breast cancer $\rightarrow$ \texttt{"breast"}\\[2pt]
         \textbf{4.}\; After tool calls $\rightarrow$
          ``write a concise summary of the gene list and pathway data''};

    \draw[flowline] (reads.south) -- (prompt.north);

    % === TOOL ===
    \node[toolpanel, below=0.4cm of prompt] (tool)
        {\textbf{Tool:}\;
         \texttt{omic\_analysis\_tool(disease, organ, cell\_type)}\\[5pt]
         {\scriptsize \textbf{Pipeline:}\;
          GLiNER NER $\rightarrow$
          OmniCellTOSG scRNA-seq query $\rightarrow$
          DESeq2 differential expression $\rightarrow$
          clusterProfiler KEGG/GO enrichment}\\[3pt]
         {\scriptsize \textbf{Returns:}\;
          sample counts (disease vs.\ normal), DEG rankings with
          log2FC \& FDR, volcano plots, KEGG/GO enrichment dotplots}\\[3pt]
         {\scriptsize \textbf{Resources:}\;
          OmniCellTOSG\; $\cdot$\;
          GLiNER\; $\cdot$\;
          DESeq2/R\; $\cdot$\;
          clusterProfiler\; $\cdot$\;
          BiomedBERT}};

    \draw[flowline] (prompt.south) -- (tool.north);

    % === WRITES ===
    \node[rwpanel, draw=blue!30!black, fill=stateBlue,
          below=0.4cm of tool] (writes)
        {\textbf{$\blacktriangle$ Writes to AgentState}
         {\scriptsize (via \texttt{\_extract\_structured\_data})}\\[5pt]
         $\bullet$\; \texttt{top\_genes[]} {\scriptsize (list[str])} ---
          ranked DEGs with log2FC \& FDR\quad
          {\scriptsize e.g.\ \texttt{[MAP3K15, CFAP47, IL3RA, \ldots]}}\\[2pt]
         $\bullet$\; \texttt{disease} {\scriptsize (str)} ---
          canonical disease name\quad
          {\scriptsize e.g.\ \texttt{"Alzheimer disease"}}\\[2pt]
         $\bullet$\; \texttt{cell\_type} {\scriptsize (str)} ---
          cell-type filter if specified\quad
          {\scriptsize e.g.\ \texttt{"microglial cell"}}\\[2pt]
         $\bullet$\; \texttt{pathways[]} {\scriptsize (list)} ---
          KEGG and GO enrichment results\\[2pt]
         $\bullet$\; \texttt{deg\_results} {\scriptsize (dict)} ---
          full differential-expression statistics per gene\\[2pt]
         $\bullet$\; \texttt{plot\_paths} {\scriptsize (dict)} ---
          volcano plot, KEGG dotplot, enrichment plot file paths};

    \draw[flowline] (tool.south) -- (writes.north);

    % === DOWNSTREAM CONSUMERS ===
    \node[consumerlabel, below=0.3cm of writes] (consumers)
        {$\downarrow$\; Consumed by:\;
         \textbf{BioMarkerKGAgent} (\texttt{top\_genes})\; $\cdot$\;
         \textbf{PubMedResearcher} (\texttt{top\_genes})\; $\cdot$\;
         \textbf{ScientistsAgent} (\texttt{top\_genes}, \texttt{pathways})\; $\cdot$\;
         \textbf{ReportingNode} (\texttt{plot\_paths})};

\end{tikzpicture}%
}% end resizebox
\caption{\textbf{OmicMiningAgent} --- workflow interface view.
The ordering shown here reflects a \emph{typical} plan generated by
the LLM orchestrator; in practice the planner dynamically decides
which agents to invoke, in what order, and whether to loop back after
failures or surprising results.  That said, the planning prompt
strongly encourages running this agent first to ``ground ALL
hypotheses in actual omics data.''  It reads only the user query and session identifier,
maps them to canonical disease/organ/cell-type forms via its system
prompt, and invokes a single tool (\texttt{omic\_analysis\_tool}) that
chains NER, database retrieval, differential expression, and pathway
enrichment.  The six fields it writes to \texttt{AgentState}---most
critically \texttt{top\_genes[]}---become the shared anchor that
every downstream agent reads.}
\label{fig:closeup-omic}
\end{figure}

% Complements internal tool illustrations with orchestrator-
% level data flow, prompt highlights, and state I/O.
% ============================================================
\begin{figure}[htbp]
\centering
\resizebox{\textwidth}{!}{%
\begin{tikzpicture}[node distance=0.35cm, font=\sffamily,
    stepbar/.style={rectangle, rounded corners=6pt, thick,
        minimum width=19cm, align=center, font=\sffamily,
        drop shadow={opacity=0.10}, inner sep=6pt, text width=18cm},
    rwpanel/.style={rectangle, rounded corners=5pt, thick,
        minimum width=19cm, align=left, inner sep=6pt,
        font=\sffamily\small, text width=17.5cm},
    infopanel/.style={rectangle, rounded corners=5pt, draw=gray!50, thick,
        minimum width=19cm, align=left, fill=white,
        inner sep=6pt, font=\sffamily\small, text width=17.5cm},
    toolpanel/.style={rectangle, rounded corners=5pt, draw=gray!60, thick,
        minimum width=19cm, align=left, fill=gray!4,
        inner sep=6pt, font=\sffamily\small, text width=17.5cm},
    flowline/.style={->, >=Stealth, very thick, draw=gray!40},
    consumerlabel/.style={font=\sffamily\small\itshape, text=gray!60!black,
        align=center, text width=19cm}
]
    % === STEP BAR ===
    \node[stepbar, draw=orange!50!black, fill=agentOrange] (step)
        {\iconDNA{orange!60!black}\;\;
         \textbf{\large OmicMiningAgent --- Data Foundation}
         {\normalsize\itshape\; (typically first)}\\[3pt]
         {\small Planner: \textit{``ALWAYS run FIRST to ground analysis
          in actual expression data --- establishes data foundation.''}}\\[2pt]
         {\scriptsize\sffamily\color{gray!60!black}
          Note: ordering shown is a \emph{typical} realisation;
          the LLM planner may reorder, skip, or loop back to any
          agent based on intermediate results.}};

    % === READS ===
    \node[rwpanel, draw=green!40!black, fill=agentGreen,
          below=0.4cm of step] (reads)
        {\textbf{$\blacktriangledown$ Reads from AgentState}\\[5pt]
         $\bullet$\; \texttt{query} {\scriptsize (str)} ---
          user's natural-language research question\\[2pt]
         $\bullet$\; \texttt{session\_id} {\scriptsize (str)} ---
          unique session identifier for file paths and result caching};

    \draw[flowline] (step.south) -- (reads.north);

    % === PROMPT HIGHLIGHTS ===
    \node[infopanel, below=0.4cm of reads] (prompt)
        {\textbf{System Prompt Highlights}
         {\scriptsize (from \texttt{OmicMiningAgent.system\_message})}\\[5pt]
         \textbf{1.}\; Extract \texttt{disease}, \texttt{organ},
          \texttt{cell\_type} from the user query\\[2pt]
         \textbf{2.}\; Map free-text to canonical forms:\;
          ``Alzheimer's'' $\rightarrow$ \texttt{"Alzheimer disease"},\;
          ``Lung cancer'' $\rightarrow$ \texttt{"lung adenocarcinoma"}\\[2pt]
         \textbf{3.}\; Map organs:\;
          Alzheimer's $\rightarrow$ \texttt{"brain"},\;
          Breast cancer $\rightarrow$ \texttt{"breast"}\\[2pt]
         \textbf{4.}\; After tool calls $\rightarrow$
          ``write a concise summary of the gene list and pathway data''};

    \draw[flowline] (reads.south) -- (prompt.north);

    % === TOOL ===
    \node[toolpanel, below=0.4cm of prompt] (tool)
        {\textbf{Tool:}\;
         \texttt{omic\_analysis\_tool(disease, organ, cell\_type)}\\[5pt]
         {\scriptsize \textbf{Pipeline:}\;
          GLiNER NER $\rightarrow$
          OmniCellTOSG scRNA-seq query $\rightarrow$
          DESeq2 differential expression $\rightarrow$
          clusterProfiler KEGG/GO enrichment}\\[3pt]
         {\scriptsize \textbf{Returns:}\;
          sample counts (disease vs.\ normal), DEG rankings with
          log2FC \& FDR, volcano plots, KEGG/GO enrichment dotplots}\\[3pt]
         {\scriptsize \textbf{Resources:}\;
          OmniCellTOSG\; $\cdot$\;
          GLiNER\; $\cdot$\;
          DESeq2/R\; $\cdot$\;
          clusterProfiler\; $\cdot$\;
          BiomedBERT}};

    \draw[flowline] (prompt.south) -- (tool.north);

    % === WRITES ===
    \node[rwpanel, draw=blue!30!black, fill=stateBlue,
          below=0.4cm of tool] (writes)
        {\textbf{$\blacktriangle$ Writes to AgentState}
         {\scriptsize (via \texttt{\_extract\_structured\_data})}\\[5pt]
         $\bullet$\; \texttt{top\_genes[]} {\scriptsize (list[str])} ---
          ranked DEGs with log2FC \& FDR\quad
          {\scriptsize e.g.\ \texttt{[MAP3K15, CFAP47, IL3RA, \ldots]}}\\[2pt]
         $\bullet$\; \texttt{disease} {\scriptsize (str)} ---
          canonical disease name\quad
          {\scriptsize e.g.\ \texttt{"Alzheimer disease"}}\\[2pt]
         $\bullet$\; \texttt{cell\_type} {\scriptsize (str)} ---
          cell-type filter if specified\quad
          {\scriptsize e.g.\ \texttt{"microglial cell"}}\\[2pt]
         $\bullet$\; \texttt{pathways[]} {\scriptsize (list)} ---
          KEGG and GO enrichment results\\[2pt]
         $\bullet$\; \texttt{deg\_results} {\scriptsize (dict)} ---
          full differential-expression statistics per gene\\[2pt]
         $\bullet$\; \texttt{plot\_paths} {\scriptsize (dict)} ---
          volcano plot, KEGG dotplot, enrichment plot file paths};

    \draw[flowline] (tool.south) -- (writes.north);

    % === DOWNSTREAM CONSUMERS ===
    \node[consumerlabel, below=0.3cm of writes] (consumers)
        {$\downarrow$\; Consumed by:\;
         \textbf{BioMarkerKGAgent} (\texttt{top\_genes})\; $\cdot$\;
         \textbf{PubMedResearcher} (\texttt{top\_genes})\; $\cdot$\;
         \textbf{ScientistsAgent} (\texttt{top\_genes}, \texttt{pathways})\; $\cdot$\;
         \textbf{ReportingNode} (\texttt{plot\_paths})};

\end{tikzpicture}%
}% end resizebox
\caption{\textbf{OmicMiningAgent} --- workflow interface view.
The ordering shown here reflects a \emph{typical} plan generated by
the LLM orchestrator; in practice the planner dynamically decides
which agents to invoke, in what order, and whether to loop back after
failures or surprising results.  That said, the planning prompt
strongly encourages running this agent first to ``ground ALL
hypotheses in actual omics data.''  It reads only the user query and session identifier,
maps them to canonical disease/organ/cell-type forms via its system
prompt, and invokes a single tool (\texttt{omic\_analysis\_tool}) that
chains NER, database retrieval, differential expression, and pathway
enrichment.  The six fields it writes to \texttt{AgentState}---most
critically \texttt{top\_genes[]}---become the shared anchor that
every downstream agent reads.}
\label{fig:closeup-omic}
\end{figure}

% Complements internal tool illustrations with orchestrator-
% level data flow, prompt highlights, and state I/O.
% ============================================================
\begin{figure}[htbp]
\centering
\resizebox{\textwidth}{!}{%
\begin{tikzpicture}[node distance=0.35cm, font=\sffamily,
    stepbar/.style={rectangle, rounded corners=6pt, thick,
        minimum width=19cm, align=center, font=\sffamily,
        drop shadow={opacity=0.10}, inner sep=6pt, text width=18cm},
    rwpanel/.style={rectangle, rounded corners=5pt, thick,
        minimum width=19cm, align=left, inner sep=6pt,
        font=\sffamily\small, text width=17.5cm},
    infopanel/.style={rectangle, rounded corners=5pt, draw=gray!50, thick,
        minimum width=19cm, align=left, fill=white,
        inner sep=6pt, font=\sffamily\small, text width=17.5cm},
    toolpanel/.style={rectangle, rounded corners=5pt, draw=gray!60, thick,
        minimum width=19cm, align=left, fill=gray!4,
        inner sep=6pt, font=\sffamily\small, text width=17.5cm},
    flowline/.style={->, >=Stealth, very thick, draw=gray!40},
    consumerlabel/.style={font=\sffamily\small\itshape, text=gray!60!black,
        align=center, text width=19cm}
]
    % === STEP BAR ===
    \node[stepbar, draw=orange!50!black, fill=agentOrange] (step)
        {\iconDNA{orange!60!black}\;\;
         \textbf{\large OmicMiningAgent --- Data Foundation}
         {\normalsize\itshape\; (typically first)}\\[3pt]
         {\small Planner: \textit{``ALWAYS run FIRST to ground analysis
          in actual expression data --- establishes data foundation.''}}\\[2pt]
         {\scriptsize\sffamily\color{gray!60!black}
          Note: ordering shown is a \emph{typical} realisation;
          the LLM planner may reorder, skip, or loop back to any
          agent based on intermediate results.}};

    % === READS ===
    \node[rwpanel, draw=green!40!black, fill=agentGreen,
          below=0.4cm of step] (reads)
        {\textbf{$\blacktriangledown$ Reads from AgentState}\\[5pt]
         $\bullet$\; \texttt{query} {\scriptsize (str)} ---
          user's natural-language research question\\[2pt]
         $\bullet$\; \texttt{session\_id} {\scriptsize (str)} ---
          unique session identifier for file paths and result caching};

    \draw[flowline] (step.south) -- (reads.north);

    % === PROMPT HIGHLIGHTS ===
    \node[infopanel, below=0.4cm of reads] (prompt)
        {\textbf{System Prompt Highlights}
         {\scriptsize (from \texttt{OmicMiningAgent.system\_message})}\\[5pt]
         \textbf{1.}\; Extract \texttt{disease}, \texttt{organ},
          \texttt{cell\_type} from the user query\\[2pt]
         \textbf{2.}\; Map free-text to canonical forms:\;
          ``Alzheimer's'' $\rightarrow$ \texttt{"Alzheimer disease"},\;
          ``Lung cancer'' $\rightarrow$ \texttt{"lung adenocarcinoma"}\\[2pt]
         \textbf{3.}\; Map organs:\;
          Alzheimer's $\rightarrow$ \texttt{"brain"},\;
          Breast cancer $\rightarrow$ \texttt{"breast"}\\[2pt]
         \textbf{4.}\; After tool calls $\rightarrow$
          ``write a concise summary of the gene list and pathway data''};

    \draw[flowline] (reads.south) -- (prompt.north);

    % === TOOL ===
    \node[toolpanel, below=0.4cm of prompt] (tool)
        {\textbf{Tool:}\;
         \texttt{omic\_analysis\_tool(disease, organ, cell\_type)}\\[5pt]
         {\scriptsize \textbf{Pipeline:}\;
          GLiNER NER $\rightarrow$
          OmniCellTOSG scRNA-seq query $\rightarrow$
          DESeq2 differential expression $\rightarrow$
          clusterProfiler KEGG/GO enrichment}\\[3pt]
         {\scriptsize \textbf{Returns:}\;
          sample counts (disease vs.\ normal), DEG rankings with
          log2FC \& FDR, volcano plots, KEGG/GO enrichment dotplots}\\[3pt]
         {\scriptsize \textbf{Resources:}\;
          OmniCellTOSG\; $\cdot$\;
          GLiNER\; $\cdot$\;
          DESeq2/R\; $\cdot$\;
          clusterProfiler\; $\cdot$\;
          BiomedBERT}};

    \draw[flowline] (prompt.south) -- (tool.north);

    % === WRITES ===
    \node[rwpanel, draw=blue!30!black, fill=stateBlue,
          below=0.4cm of tool] (writes)
        {\textbf{$\blacktriangle$ Writes to AgentState}
         {\scriptsize (via \texttt{\_extract\_structured\_data})}\\[5pt]
         $\bullet$\; \texttt{top\_genes[]} {\scriptsize (list[str])} ---
          ranked DEGs with log2FC \& FDR\quad
          {\scriptsize e.g.\ \texttt{[MAP3K15, CFAP47, IL3RA, \ldots]}}\\[2pt]
         $\bullet$\; \texttt{disease} {\scriptsize (str)} ---
          canonical disease name\quad
          {\scriptsize e.g.\ \texttt{"Alzheimer disease"}}\\[2pt]
         $\bullet$\; \texttt{cell\_type} {\scriptsize (str)} ---
          cell-type filter if specified\quad
          {\scriptsize e.g.\ \texttt{"microglial cell"}}\\[2pt]
         $\bullet$\; \texttt{pathways[]} {\scriptsize (list)} ---
          KEGG and GO enrichment results\\[2pt]
         $\bullet$\; \texttt{deg\_results} {\scriptsize (dict)} ---
          full differential-expression statistics per gene\\[2pt]
         $\bullet$\; \texttt{plot\_paths} {\scriptsize (dict)} ---
          volcano plot, KEGG dotplot, enrichment plot file paths};

    \draw[flowline] (tool.south) -- (writes.north);

    % === DOWNSTREAM CONSUMERS ===
    \node[consumerlabel, below=0.3cm of writes] (consumers)
        {$\downarrow$\; Consumed by:\;
         \textbf{BioMarkerKGAgent} (\texttt{top\_genes})\; $\cdot$\;
         \textbf{PubMedResearcher} (\texttt{top\_genes})\; $\cdot$\;
         \textbf{ScientistsAgent} (\texttt{top\_genes}, \texttt{pathways})\; $\cdot$\;
         \textbf{ReportingNode} (\texttt{plot\_paths})};

\end{tikzpicture}%
}% end resizebox
\caption{\textbf{OmicMiningAgent} --- workflow interface view.
The ordering shown here reflects a \emph{typical} plan generated by
the LLM orchestrator; in practice the planner dynamically decides
which agents to invoke, in what order, and whether to loop back after
failures or surprising results.  That said, the planning prompt
strongly encourages running this agent first to ``ground ALL
hypotheses in actual omics data.''  It reads only the user query and session identifier,
maps them to canonical disease/organ/cell-type forms via its system
prompt, and invokes a single tool (\texttt{omic\_analysis\_tool}) that
chains NER, database retrieval, differential expression, and pathway
enrichment.  The six fields it writes to \texttt{AgentState}---most
critically \texttt{top\_genes[]}---become the shared anchor that
every downstream agent reads.}
\label{fig:closeup-omic}
\end{figure}
